\documentclass[10pt]{article}

\usepackage{parskip}
\usepackage[margin=1in]{geometry} 
\usepackage{amsmath,amsthm,amssymb, graphicx, multicol, array}
\usepackage{enumitem}
\usepackage{amssymb}
\usepackage{float}
\usepackage{href-ul}
 
\newcommand{\N}{\mathbb{N}}
\newcommand{\Z}{\mathbb{Z}}
 
\newenvironment{problem}[2][Problem]{\begin{trivlist}
\item[\hskip \labelsep {\bfseries #1}\hskip \labelsep {\bfseries #2.}]}{\end{trivlist}}

\begin{document}
 
\title{Introduction}
\author{Jakob Sverre Alexandersen\\
GRA4158 Marketing and the Analysis of Experiments and Quasi-Experiments\\
Lecture 1}
\maketitle

\tableofcontents
\newpage
\section{General content overview}
\begin{itemize}
    \item What is marketing?
    \begin{itemize}
        \item Why and when is data science applied in marketing?
    \end{itemize}
    \item Foundations of causal inference
    \begin{itemize}
        \item What is causality? Differences between causal, predictive and descriptive questions 
        \item Rubin's potential outcome framework, different types of treatment effects 
        \item Randomization, selection effects, and basic inference 
        \item Basics of causal diagrams (DAGs), good and bad controls 
    \end{itemize}
    \item Experimental designs 
    \begin{itemize}
        \item Field vs. lab vs natural experiments, between vs. within subject design, multi-factorial design 
        \item internal and expernal validity of different designs
    \end{itemize}
    \item Quasi-experimental and observational methods for causal inference
    \begin{itemize}
        \item regression adjustment, matching and weighing, causal ML (double debiased ML), instrumental variables (IVs) and IV-free methods, difference-in-difference, synthetic control, and synthetic difference-in-difference, regression discontinuity design
    \end{itemize}
\end{itemize}

\section{Causal questions}
\subsection{What is causality?}
\begin{itemize}
    \item In this course, we are not concerned about the nature of causation or philosophical foundations of causality 
    \item But, how to measure the effect of a cause? 
    \begin{itemize}
        \item To measure causal effects, we need a presumed effect (Y) and a presumed cause (X)
    \end{itemize}
    \item In social sciences (including business and marketing), causal questions typically matter more than associational questions 
    \item We usually do not only want to learn, if two concepts are related 
    \item We want to learn if changing one aspect leads to an outcome
    \begin{itemize}
        \item Inform managerial or policy decisions
    \end{itemize}
\end{itemize}
\textbf{Examples}
\begin{itemize}
    \item Descriptive (correlational) question: Do countries with higher minimum wages have less poverty?
    \item Causal question: does raising the minimum wage reduce poverty?
    \item Descriprive (correlational) question: do people who sleep less die early?
    \item Causal question: does a lack of sleep reduce the (expected) lifetime?
\end{itemize}
\newpage 
\subsection{Causal questions in marketing}
\begin{itemize}
    \item we do not only want to know... 
    \begin{itemize}
        \item if people who saw an advertisement buy more rp
        \item This is all just a bunch of nonsense – in the regard that this is all very obvious and self-insisting. I will only take notes on stuff that I do not know 
    \end{itemize}
\end{itemize}
\section{Another view on causality}
\begin{itemize}
    \item Three classic conditions for causality from Kenny (1979):
    \begin{enumerate}
        \item X must precede Y temporally 
        \item X must be reliably correlated with Y (beyond chance)
        \item The relation between X and Y must not be spurious (based on false reasoning or information that is not true)
    \end{enumerate}
\end{itemize}
\subsection{The three classic conditions}
\subsubsection{X must precede Y temporally}
\begin{itemize}
    \item Changes in X must happen before we observe the changes in Y 
    \item A typical criticism about condition 1 is that it ignores the possibility of ``simultaneity''
    \begin{itemize}
        \item philosophical debate about whether instantaneous causation is possible 
        \item reciprocal relationship (feedback-loops)
        \begin{itemize}
            \item temporal precedence does not forbid that Y causes X in return, but just at a different point in time 
            \item $X_{t1} \to X_{t2} \to X_{t3} ... $
        \end{itemize}
        \item simultaneity is often just a problem when the available temporal resolution of our measurement is not large enough to distinguish cause and effect 
    \end{itemize}
    \item it is also one of the most common fallacies to think that close temporal proximity implies causation
\end{itemize}
\newpage 
\subsubsection{X must be reliably correlated with Y (beyond chance)}
\begin{itemize}
    \item generally, causation requires some co-variation (correlation) or co-occurence of events
    \begin{itemize}
        \item if we observe A, we also observe B
        \item when the cause changes, we can also observe a change in the effect 
        \item when we press the light button, the light goes off 
        \item when we reduce advertising budget, sales go down 
        \item when we implement a loyalty program, sales increase 
    \end{itemize}
    \item however, co-occurence is neither necessary nor sufficient 
    \item just because people who drown and the number of movies Nic Cage is starred in are correlated, does not mean that Nic Cage drowns people
    \item Another problem: X can cause Y, even in the absence of any observable correlation in the data 
    \begin{itemize}
        \item ... in the absence of an observable main effect
        \item e.g., effect heterogeneity (moderation); competitive mediation; non-linear relations
    \end{itemize}
    \item Kenny (2004) speaks more generally about ``the presence of a functional relationship between cause and effect''
    \begin{itemize}
        \item ``two variables are independent if knowing the value of one variable provides no information about the value of the other variable''
        \item functional relationship can be more complex that linear / correlational
    \end{itemize}
\end{itemize}
\subsubsection{The relation between X and Y must not be spurious}
\begin{itemize}
    \item The relation is not explained by other causes 
    \begin{itemize}
        \item typical correlation vs. causation problem 
    \end{itemize}
    \item this will be one of our main concerns during this course 
    \begin{itemize}
        \item how can we achieve that the relationship between X and Y is not confounded by other variables?
    \end{itemize}
\end{itemize}
\subsection{Confounder}
\begin{itemize}
    \item Z is a confounder to the X-Y relationship if it influences both X and Y 
    \begin{itemize}
        \item Z is a common cause of both X and Y 
    \end{itemize}
    \item an observed correlation between X and Y partly (or only) exists because of Z
\end{itemize}
\newpage 
\section{Causal vs. non-causal wording}
\begin{itemize}
    \item There are lots of words that are generally taken to imply causality, as well as words that describe relationships without implying causality.
    \item A causal statement has two components: a cause and and effect 
    \item Causal wording 
    \begin{itemize}
        \item we can replace X ``causes'' Y by... 

        affects, increases, decreases, changes, leads to, improves
    \end{itemize}
    \item non causal wording (i.e., correlational), but often falsely understood to be causal 
    \begin{itemize}
        \item X and Y are associated, related, tend to occur together, X is linked to Y, X predicts Y, etc
    \end{itemize}
\end{itemize}
\section{Causal vs. predictive}
\begin{itemize}
    \item In this course, we want to understand whether A causes B 
    \item the questions whether A can predict B is somthing different 
    \item descriptive question: do countries with higher minimum wage have less poverty?
    \item causal question: does 
\end{itemize}

\end{document}